\chapter{Background}
\label{ch:background}

In this chapter, I provide the fundamental information 
regarding the objects considered for the subsequent analyses presented in this thesis: tarballs.

An \ac{NPM} tarball, according to the official \ac{NPM} documentation \cite{npmdocs2026npmdocumentation}, 
is a compressed archive with a .tgz extension, 
containing all the files needed to install an \ac{NPM} package.


Although the tarball and the source repository of an \ac{NPM} package are related concepts, they are distinct.
In fact, the source repository represents the development process, while the \ac{NPM} tarball represents the distributed product.

To ensure a fair and symmetrical comparison, it is necessary to analyze the same type of artifact,
as source repositories and published tarballs can differ significantly in the files they include.

According to the official \ac{NPM} documentation \cite{npmdocs2026npmdocumentation},
and as confirmed by the analysis I conducted in \nameref{ch:analysis},
an \ac{NPM} tarball generally contains the following categories of files:

\begin{itemize}
    \item \textbf{package.json:}
    It is mandatory to have this file present inside the tarball.
    Publishing a package requires that the \texttt{package.\allowbreak json} file must at least specify the fields for the package name and its version.
    Other common optional fields are the description, license, repository URL, any dependencies
    from other packages, and more.
    Indeed, it is not necessary for an \ac{NPM} package to declare a link to a public repository.


    \item \textbf{Code files:}
    The tarball generally includes JavaScript and/or TypeScript code files.
    The code is distributed and included directly in the tarball itself, for reasons of transparency, inspectability, and debugging.
    However, as \textcite{10.5555/3361338.3361407} highlights,
    there is no guarantee that the code present in the \ac{NPM} tarball is equivalent to the one in the public version control system (e.g. GitHub).
    
    Furthermore, the code can be manually altered after its 
    publication, for example by an attacker injecting malicious code;
    and the user cannot verify this with certainty.
    The presence and structure of these files vary by package, as the standard is not rigid.

    Rarely, some packages may not have any JavaScript or TypeScript code files at all, 
    and this depends on the purpose of the package.
    For example, there are packages whose resource and utility lie solely in the data inside \texttt{.json} files.
    Notable examples include \texttt{node-\allowbreak releases}, which offers a list of Node.js release versions, 
    and \texttt{spdx-\allowbreak license-\allowbreak ids}, which provides the official list of \ac{SPDX} licenses. 
    Developers use these packages as dependencies to get the current data without hard-coding it manually.

    Alternatively, some packages may have no code at all because they are deprecated and maintained solely as placeholders, 
    like recently \texttt{@types/\allowbreak uuid} package.


    \item \textbf{Metadata:}
    Frequently, documentation files, 
    such as Markdown files like \texttt{README.md} and \texttt{CHANGELOG.md}, and the \texttt{LICENSE} text file 
    used to specify copyright; are included inside the tarball.

    \item \textbf{Build files/artifacts:}
    Spesso, all'interno del tarball npm, si trovano file che non sono scritti direttamente dagli sviluppatori, 
    ma vengono generati tramite un processo di build o compilazione a partire dal codice sorgente originale.

    Questi artefatti permettono agli utenti di integrare e utilizzare i pacchetti nei propri progetti senza 
    dover replicare il processo di build originale, che potrebbe richiedere configurazioni complesse di tool come Webpack, Rollup o TypeScript. 
    Ciò garantisce che le funzionalità del pacchetto siano immediatamente fruibili in tutti gli ambienti di 
    destinazione, come in Node.js, browser, ecc.

    In alcuni casi, per ridurre le dimensioni del pacchetto, 
    gli sviluppatori possono scegliere di escludere completamente il codice sorgente, e includere solo i file compilati,
    come nel package @babel/parser dalla versione 8.0.0-beta.2 alla versione 8.0.0-beta.3.
    Qui inoltre si può vedere che questi artifacts are not necessarily \texttt{.js} files, ma 
    they can also be source maps (\texttt{.map}) or declaration files (\texttt{.d.ts}).

    Non tutti i pacchetti su npm contengono artefatti di build; infatti, 
    se il codice è scritto in JavaScript conforme agli standard supportati (ECMAScript) 
    e non utilizza linguaggi che richiedono traspilazione (come TypeScript), 
    può essere distribuito direttamente nella sua forma originale.

    These artifacts possono variare da pacchetto a pacchetto e tra versioni diverse dello stesso pacchetto.
    Ad esempio possono diventare più ottimizzati tra una versione e l'altra; for example, moving from multi-line files with spaces and comments
    to a single line with no comments and minimal whitespace; as in the case of the \texttt{svgo} package for the file \texttt{dist/svgo.browser.js}
    from version 3.3.1 to 3.3.2.
    Oppure rimanere uguale, ad esempio the \texttt{css-tree} package for the file \texttt{dist/csstree.js} che sta su più righe. 
    
    Inoltre, as shown by \textcite{9240695}, even if the source code remains the same, the build files, and 
    consequently the final tarball, may differ due to non-deterministic factors in the build process.

    Since the \ac{NPM} ecosystem is permissive, there is no universal way to distinguish if a file was built or hand-written simply by looking 
    at the directory or the metadata in \texttt{package.json}. 
    As highlighted by \textcite{9240695}, common conventions suggest that folders such as \texttt{dist/}, \texttt{build/} or \texttt{lib/} contain artifacts, 
    but this division is not guaranteed or verifiable a priori, because it depends on the developer's choices.
    
    In fact, in Chapter \ref{ch:analysis}, we will see that generated artifacts can reside in other folders; 
    for example the \texttt{css-what} package contains the generated \texttt{out.json} file in \texttt{src/( \_\_fixtures\_\_/)}.
    Additionally, there are cases where a hand-written files reside in \texttt{dist/} or \texttt{build/}, 
    such as the \texttt{helper\allowbreak .d.ts} file for the \texttt{domutils} package, which is located in the \texttt{lib/} directory;
    or the \texttt{lib/\allowbreak index\allowbreak .d.ts} file for the \texttt{@babel/\allowbreak parser} package from version 8.0.0-beta.2 to 8.0.0-beta.3.
    
    Hence, I decided not to exclude any folders based on their name, avoiding the risk of overlooking relevant files.
    

\end{itemize}

As can be observed better in \nameref{ch:analysis}, other files can be found in tarballs.

The \texttt{html} files, after careful analysis,
represent mainly test files, code coverage reports, and demos, such in the \texttt{jwt-decode} package;
or presentation pages of the package, documentation, and licenses, such in the \texttt{sax} and \texttt{pino} packages.

Instead, the \texttt{css} files are either used as test fixtures (e.g., example data for testing), 
as in the case of the \texttt{convert-source-map} package;
or simply as styles for HTML-based tests and demos, as in the \texttt{loglevel} package.

However, in modern packaging practices,
non essential assets are often removed to keep the final product as lightweight as possible.

Moreover, like build artifacts, these files are not always located in specific directories,
and it is not possible to determine with certainty if a file is test file or not.




\begin{comment}
Third-party packages play a crucial role in modern software
development by reducing implementation effort and providing reusable functionalities 
that would otherwise be timeconsuming to develop from scratch \cite{unknown}.

To support JavaScript developers in the reuse of third-party code, the \ac{NPM} provides
hundreds of thousands of free and reusable code software packages.
The \ac{NPM} platform consists of an online registry for searching 
packages suitable for given tasks, and a package manager
that automatically resolves and installs dependencies, 
simplify the integration of these packages into projects \cite{10.5555/3361338.3361407}.
By design, The \ac{NPM} ecosystem is open, allowing any user to freely publish.

This openness has been a key factor in the rapid growth and success of \ac{NPM},
but it also introduces significant security risks.
In particular, the widespread use of third-party dependencies makes \ac{NPM}
a central component of the JavaScript software supply chain.

The software Supply Chain \ac{SSC} refers to the set of processes to select and obtain software components
from third parties; it also comprehends the companies involved in these processes \cite{gokkaya2023softwaresupplychainreview}.

Modern software systems heavily rely on \ac{SSC}, often incorporating a large number of direct
and transitive dependencies.
Within this context, the \ac{NPM} ecosystem exemplifies how supply chain dependencies
can amplify the impact of security incidents.

Highly popular packages can directly or indirectly influence
many other projects, sometimes exceeding 100,000 dependent packages \cite{10.5555/3361338.3361407}.
Moreover, installing an average \ac{NPM} package implicitly introduces
trust on 79 third-party packages and 39 maintainers \cite{10.5555/3361338.3361407}.

As a consequence, the \ac{NPM} ecosystem has become a prominent target for
Software Supply Chain attacks \acp{SCCA}, characterized by the injection of malicious code
into legitimate packages update, in order to compromise
dependent systems further down the chain \cite{ohm2020backstabbersknifecollectionreview}.
\end{comment}

\begin{comment}

This chapter provides the necessary background information to understand the content of your thesis.
The goal is to make your thesis \emph{self-contained}: the reader, who you can assume will be
a fellow Computer Science student, should be able to understand your work without having to refer
to external sources. Of course, you should reference the sources you used to write this chapter, and references to documents that the reader can use to know more about those topics.

Do not include information that is not necessary to understand your work. If something is complementary to your work, but not useful to understand it, you can include it in the \nameref{ch:related} (\Cref{ch:related}).

Guide the reader while they are going through this chapter: do not just dump information, but present the context in which they will be used in the thesis, and what those concepts will be useful for.\footnote{A good exercise is thinking that your reader will constantly ask ``so what?'' after everything you tell them. If it is unclear why you provided them with information, you should explain (or, of course, remove the content if there is no explanation).}
Put forward references to the places these concepts are used, and put backward references from that chapter to here so that the reader can link the pieces.\footnote{Always keep in mind that the reader will not necessarily read all your documents. Such forward- and back-references help refresh the memory of those who read your documents linearly \emph{and} help those who do not.}
Pay close attention, here and everywhere else in your thesis, to make sure that every non-trivial concept you use is defined clearly before being used.

Sometimes, there will be concepts that are tangential to your work: in those cases, avoid having an unclear or incomplete explanation: either explain it clearly in sufficient depth, or skip the concept altogether, leaving a bibliographic reference for those who want to understand the issue. A short but unclear explanation that will confuse readers is worse than either choice.

The Background chapter contains information about work that \emph{you didn't do}. The work you
did should be described in the \nameref{ch:method} (\Cref{ch:method}) instead.

The amount of information you will need to present in this chapter can vary a lot between theses.
In some cases, you will have no or very little need for a Background chapter; in that case, it
may make sense to move this information to a section in the Related Work chapter.
\end{comment}